% GENERATED FILE. Edit canonical lessons, then run npm run generate:books.
% canonical-source-hash: fnv1a64:3bab2d4db50a4768
% canonical-lessons: TE-C461-attapette, TE-C461-vadilivellu, TE-C461-paresuko, TE-C461-taralincu, TE-C461-ceraveyu, TE-R461-review-454-461

\chapter{Bags and Boxes}
\label{ch:te-bags-and-boxes}
\hlchaptermodality{461}

\begin{chapteropening}
\textbf{By the end of this chapter:} \emph{I can say a cardboard box, to leave behind, to mislay, to shift elsewhere and to deliver.}

Complete the last lesson of chapter 461: I can say a cardboard box, to leave behind, to mislay, to shift elsewhere and to deliver.
\end{chapteropening}

\section[\texorpdfstring{\te{అట్టపెట్టె} (aṭṭapeṭṭe)}{అట్టపెట్టె (aṭṭapeṭṭe)}]{\te{అట్టపెట్టె} (aṭṭapeṭṭe) — a cardboard box}

\label{lesson:TE-C461-attapette}



\begin{quote}
Before the new one: say the Telugu for a refund, then the Telugu for to double-check.
\end{quote}

\subsection*{You'll want to know: \te{అట్టపెట్టె}}
\textbf{\te{అట్టపెట్టె}} — \emph{aṭṭapeṭṭe} — \textquotedblleft{}a cardboard box\textquotedblright{}.

\begin{quote}
\te{పుస్తకాలను} \te{ఒక} \te{అట్టపెట్టెలో} \te{పెట్టాం}.
\end{quote}

\begin{quote}
\emph{pustakālanu oka aṭṭapeṭṭelō peṭṭāṁ.}
\end{quote}

\begin{quote}
We put the books in a cardboard box.
\end{quote}

\te{అట్ట}, cardboard, + \te{పెట్టె}, a box.

\begin{cousinweb}
\textbf{\te{అట్టపెట్టె}} (\emph{aṭṭapeṭṭe}) is native Dravidian: \te{అట్ట}, \textquotedblleft{}leather, cardboard\textquotedblright{} (DEDR 98), and \te{పెట్టె}, \textquotedblleft{}a box\textquotedblright{} (DEDR 4388).
\end{cousinweb}

\begin{grammarlens}[title={What you've built}]
One more word for getting around.
\end{grammarlens}

\subsection*{Your turn}
\begin{itemize}
\raggedright
  \item \emph{Say it:} \emph{aṭṭapeṭṭe}
  \item \emph{Say it:} \emph{aṭṭapeṭṭe}, once more
  \item \emph{Say it:} \emph{pustakālanu oka aṭṭapeṭṭelō peṭṭāṁ.} — We put the books in a cardboard box.
  \item \emph{Recall:} say the Telugu for a refund, then the Telugu for to double-check, then say \emph{aṭṭapeṭṭe} again
\end{itemize}

\begin{tcolorbox}[breakable,colback=teal!4,colframe=teal!35!black,title={Before you move on}]
What does \te{అట్టపెట్టె} mean? (\textquotedblleft{}A cardboard box\textquotedblright{}.) Say it once more.
\end{tcolorbox}
\section[\texorpdfstring{\te{వదిలివెళ్ళు} (vadiliveḷḷu)}{వదిలివెళ్ళు (vadiliveḷḷu)}]{\te{వదిలివెళ్ళు} (vadiliveḷḷu) — to leave behind}

\label{lesson:TE-C461-vadilivellu}



\begin{quote}
Before the new one: say the Telugu for to double-check, then the Telugu for a cardboard box.
\end{quote}

\subsection*{You'll want to know: \te{వదిలివెళ్ళు}}
\textbf{\te{వదిలివెళ్ళు}} — \emph{vadiliveḷḷu} — \textquotedblleft{}to leave behind\textquotedblright{}.

\begin{quote}
\te{నా} \te{గొడుగు} \te{బస్సులో} \te{వదిలివెళ్ళాను}.
\end{quote}

\begin{quote}
\emph{nā goḍugu bassulō vadiliveḷḷānu.}
\end{quote}

\begin{quote}
I left my umbrella behind on the bus.
\end{quote}

Past (I …): \te{వదిలివెళ్ళాను}, \emph{vadiliveḷḷānu}.

\te{వదిలి}, having left, + \te{వెళ్ళు}, to go: to go and leave something.

\begin{cousinweb}
\textbf{\te{వదిలివెళ్ళు}} (\emph{vadiliveḷḷu}) is native Dravidian: \te{వదిలి}, from \te{వదులు}, \textquotedblleft{}to let go, to leave\textquotedblright{} (DEDR 5243), and \te{వెళ్ళు}, \textquotedblleft{}to go\textquotedblright{}.
\end{cousinweb}

\begin{grammarlens}[title={What you've built}]
One more word for getting around.
\end{grammarlens}

\subsection*{Your turn}
\begin{itemize}
\raggedright
  \item \emph{Say it:} \emph{vadiliveḷḷu}
  \item \emph{Say it:} \emph{vadiliveḷḷu}, once more
  \item \emph{Say it:} \emph{nā goḍugu bassulō vadiliveḷḷānu.} — I left my umbrella behind on the bus.
  \item \emph{Recall:} say the Telugu for to double-check, then the Telugu for a cardboard box, then say \emph{vadiliveḷḷu} again
\end{itemize}

\begin{tcolorbox}[breakable,colback=teal!4,colframe=teal!35!black,title={Before you move on}]
What does \te{వదిలివెళ్ళు} mean? (\textquotedblleft{}To leave behind\textquotedblright{}.) Say it once more.
\end{tcolorbox}
\section[\texorpdfstring{\te{పారేసుకో} (pārēsukō)}{పారేసుకో (pārēsukō)}]{\te{పారేసుకో} (pārēsukō) — to mislay; to lose by carelessness}

\label{lesson:TE-C461-paresuko}



\begin{quote}
Before the new one: say the Telugu for a cardboard box, then the Telugu for to leave behind.
\end{quote}

\subsection*{You'll want to know: \te{పారేసుకో}}
\textbf{\te{పారేసుకో}} — \emph{pārēsukō} — \textquotedblleft{}to mislay; to lose by carelessness\textquotedblright{}.

\begin{quote}
\te{అతను} \te{స్టేషన్లో} \te{పర్సు} \te{పారేసుకున్నాడు}.
\end{quote}

\begin{quote}
\emph{atanu sṭēṣanlō parsu pārēsukunnāḍu.}
\end{quote}

\begin{quote}
He lost his purse at the station.
\end{quote}

Past (I …): \te{పారేసుకున్నాను}, \emph{pārēsukunnānu}.

\te{పారేయు}, to throw away, + -\te{కో}: to lose something of one's own by carelessness.

\begin{cousinweb}
\textbf{\te{పారేసుకో}} (\emph{pārēsukō}) is a Telugu verb built on \te{పారేయు}, \textquotedblleft{}to throw away\textquotedblright{}, a compound whose first part has an origin that is not settled and whose second part is \te{వేయు}, \textquotedblleft{}to put\textquotedblright{} (DEDR 5549), with -\te{కో}, \textquotedblleft{}for oneself\textquotedblright{} (DEDR 2151).
\end{cousinweb}

\begin{grammarlens}[title={What you've built}]
One more word for getting around.
\end{grammarlens}

\subsection*{Your turn}
\begin{itemize}
\raggedright
  \item \emph{Say it:} \emph{pārēsukō}
  \item \emph{Say it:} \emph{pārēsukō}, once more
  \item \emph{Say it:} \emph{atanu sṭēṣanlō parsu pārēsukunnāḍu.} — He lost his purse at the station.
  \item \emph{Recall:} say the Telugu for a cardboard box, then the Telugu for to leave behind, then say \emph{pārēsukō} again
\end{itemize}

\begin{tcolorbox}[breakable,colback=teal!4,colframe=teal!35!black,title={Before you move on}]
What does \te{పారేసుకో} mean? (\textquotedblleft{}To mislay; to lose by carelessness\textquotedblright{}.) Say it once more.
\end{tcolorbox}
\section[\texorpdfstring{\te{తరలించు} (taraliñcu)}{తరలించు (taraliñcu)}]{\te{తరలించు} (taraliñcu) — to shift elsewhere; to move}

\label{lesson:TE-C461-taralincu}



\begin{quote}
Before the new one: say the Telugu for to leave behind, then the Telugu for to mislay.
\end{quote}

\subsection*{You'll want to know: \te{తరలించు}}
\textbf{\te{తరలించు}} — \emph{taraliñcu} — \textquotedblleft{}to shift elsewhere; to move\textquotedblright{}.

\begin{quote}
\te{వరద} \te{ముందు} \te{ప్రజలను} \te{సురక్షిత} \te{ప్రాంతాలకు} \te{తరలించారు}.
\end{quote}

\begin{quote}
\emph{varada mundu prajalanu surakṣita prāntālaku taraliñcāru.}
\end{quote}

\begin{quote}
Before the flood people were moved to safe places.
\end{quote}

Past (I …): \te{తరలించాను}, \emph{taraliñcānu}.

Said of moving people, goods or an office to another place.

\begin{cousinweb}
\textbf{\te{తరలించు}} (\emph{taraliñcu}) is native Dravidian: DEDR 3128 lists Telugu \emph{taralu}, \textquotedblleft{}to set out, to start\textquotedblright{}, here made causative with -\te{ఇంచు}: to make something move off.
\end{cousinweb}

\begin{grammarlens}[title={What you've built}]
One more word for getting around.
\end{grammarlens}

\subsection*{Your turn}
\begin{itemize}
\raggedright
  \item \emph{Say it:} \emph{taraliñcu}
  \item \emph{Say it:} \emph{taraliñcu}, once more
  \item \emph{Say it:} \emph{varada mundu prajalanu surakṣita prāntālaku taraliñcāru.} — Before the flood people were moved to safe places.
  \item \emph{Recall:} say the Telugu for to leave behind, then the Telugu for to mislay, then say \emph{taraliñcu} again
\end{itemize}

\begin{tcolorbox}[breakable,colback=teal!4,colframe=teal!35!black,title={Before you move on}]
What does \te{తరలించు} mean? (\textquotedblleft{}To shift elsewhere; to move\textquotedblright{}.) Say it once more.
\end{tcolorbox}
\section[\texorpdfstring{\te{చేరవేయు} (cēravēyu)}{చేరవేయు (cēravēyu)}]{\te{చేరవేయు} (cēravēyu) — to deliver; to carry to}

\label{lesson:TE-C461-ceraveyu}



\begin{quote}
Before the new one: say the Telugu for to mislay, then the Telugu for to shift elsewhere.
\end{quote}

\subsection*{You'll want to know: \te{చేరవేయు}}
\textbf{\te{చేరవేయు}} — \emph{cēravēyu} — \textquotedblleft{}to deliver; to carry to\textquotedblright{}.

\begin{quote}
\te{కూలీలు} \te{సామాను} \te{రైలు} \te{దగ్గరికి} \te{చేరవేశారు}.
\end{quote}

\begin{quote}
\emph{kūlīlu sāmānu railu daggariki cēravēśāru.}
\end{quote}

\begin{quote}
The porters carried the luggage to the train.
\end{quote}

Past (I …): \te{చేరవేశాను}, \emph{cēravēśānu}.

\te{చేర}, to the place, + \te{వేయు}, to put: to get something to where it must go.

\begin{cousinweb}
\textbf{\te{చేరవేయు}} (\emph{cēravēyu}) is native Dravidian: \te{చేర}, from \te{చేరు}, \textquotedblleft{}to reach\textquotedblright{} (DEDR 2814), and \te{వేయు}, \textquotedblleft{}to put\textquotedblright{} (DEDR 5549).
\end{cousinweb}

\begin{grammarlens}[title={What you've built}]
One more word for getting around.
\end{grammarlens}

\subsection*{Your turn}
\begin{itemize}
\raggedright
  \item \emph{Say it:} \emph{cēravēyu}
  \item \emph{Say it:} \emph{cēravēyu}, once more
  \item \emph{Say it:} \emph{kūlīlu sāmānu railu daggariki cēravēśāru.} — The porters carried the luggage to the train.
  \item \emph{Recall:} say the Telugu for to mislay, then the Telugu for to shift elsewhere, then say \emph{cēravēyu} again
\end{itemize}

\begin{tcolorbox}[breakable,colback=teal!4,colframe=teal!35!black,title={Before you move on}]
What does \te{చేరవేయు} mean? (\textquotedblleft{}To deliver; to carry to\textquotedblright{}.) Say it once more.
\end{tcolorbox}
\section[(dialogue)]{Words from chapters 454–461 — a review}

\label{lesson:TE-R461-review-454-461}



\begin{quote}
Say the Telugu for a pilot and international.
\end{quote}

\subsection*{Your turn}
Say these aloud:

\begin{itemize}
\raggedright
  \item a pilot, international, one's homeland, transport, to see off
  \item a tourist, a tour, a tour guide, entry, a pleasure trip
  \item a jaunt, to stroll, a view, dusk, a stopover
  \item lodging, an inn, a reception desk, an elevator, to stay overnight
  \item to settle in, to get used to, a mosquito net, a toilet, drinking water
  \item a snack, a dining hall, a facility, the web, a phone charger
  \item inconvenience, trouble, to inquire, a refund, to double-check
  \item a cardboard box, to leave behind, to mislay, to shift elsewhere, to deliver
\end{itemize}

\begin{tcolorbox}[breakable,colback=teal!4,colframe=teal!35!black,title={Before you move on}]
A pilot? (\textbf{\te{పైలటు}}, \emph{pailaṭu}.) International? (\textbf{\te{అంతర్జాతీయ}}, \emph{antarjātīya}.) One's homeland? (\textbf{\te{స్వదేశం}}, \emph{svadēśaṁ}.) Transport? (\textbf{\te{రవాణా}}, \emph{ravāṇā}.) To see off? (\textbf{\te{సాగనంపు}}, \emph{sāganampu}.) A tourist? (\textbf{\te{పర్యాటకుడు}}, \emph{paryāṭakuḍu}.) A tour? (\textbf{\te{పర్యటన}}, \emph{paryaṭana}.) A tour guide? (\textbf{\te{గైడు}}, \emph{gaiḍu}.) Entry? (\textbf{\te{ప్రవేశం}}, \emph{pravēśaṁ}.) A pleasure trip? (\textbf{\te{విహారయాత్ర}}, \emph{vihārayātra}.) A jaunt? (\textbf{\te{షికారు}}, \emph{ṣikāru}.) To stroll? (\textbf{\te{విహరించు}}, \emph{vihariñcu}.) A view? (\textbf{\te{దృశ్యం}}, \emph{dṛśyaṁ}.) Dusk? (\textbf{\te{సంధ్య}}, \emph{sandhya}.) A stopover? (\textbf{\te{మజిలీ}}, \emph{majilī}.) Lodging? (\textbf{\te{వసతి}}, \emph{vasati}.) An inn? (\textbf{\te{సత్రం}}, \emph{satraṁ}.) A reception desk? (\textbf{\te{రిసెప్షను}}, \emph{risepṣanu}.) An elevator? (\textbf{\te{లిఫ్టు}}, \emph{liphṭu}.) To stay overnight? (\textbf{\te{బస} \te{చేయు}}, \emph{basa cēyu}.) To settle in? (\textbf{\te{సర్దుకో}}, \emph{sardukō}.) To get used to? (\textbf{\te{అలవాటుపడు}}, \emph{alavāṭupaḍu}.) A mosquito net? (\textbf{\te{దోమతెర}}, \emph{dōmatera}.) A toilet? (\textbf{\te{మరుగుదొడ్డి}}, \emph{marugudoḍḍi}.) Drinking water? (\textbf{\te{తాగునీరు}}, \emph{tāgunīru}.) A snack? (\textbf{\te{ఫలహారం}}, \emph{phalahāraṁ}.) A dining hall? (\textbf{\te{భోజనశాల}}, \emph{bhōjanaśāla}.) A facility? (\textbf{\te{సదుపాయం}}, \emph{sadupāyaṁ}.) The web? (\textbf{\te{అంతర్జాలం}}, \emph{antarjālaṁ}.) A phone charger? (\textbf{\te{ఛార్జరు}}, \emph{chārjaru}.) Inconvenience? (\textbf{\te{అసౌకర్యం}}, \emph{asaukaryaṁ}.) Trouble? (\textbf{\te{ఇబ్బంది}}, \emph{ibbandi}.) To inquire? (\textbf{\te{వాకబు} \te{చేయు}}, \emph{vākabu cēyu}.) A refund? (\textbf{\te{వాపసు}}, \emph{vāpasu}.) To double-check? (\textbf{\te{సరిచూసుకో}}, \emph{saricūsukō}.) A cardboard box? (\textbf{\te{అట్టపెట్టె}}, \emph{aṭṭapeṭṭe}.) To leave behind? (\textbf{\te{వదిలివెళ్ళు}}, \emph{vadiliveḷḷu}.) To mislay? (\textbf{\te{పారేసుకో}}, \emph{pārēsukō}.) To shift elsewhere? (\textbf{\te{తరలించు}}, \emph{taraliñcu}.) To deliver? (\textbf{\te{చేరవేయు}}, \emph{cēravēyu}.)
\end{tcolorbox}
